% !TeX root = ../papers/nc-p3.tex
\section{四変数全体での中心四次式とPI恒等式}
\label{sec:klein-generic-central-parameters}
\begin{theorem}\label{thm:klein-generic-central-parameters}
四次元の第二補完族の幾何学的一般点で、四頂点環の四次中心の文字次数別次元は
$(3,2,0,0)$である。定理\ref{thm:klein-rational-period-four}の有理点を含む
単一の稠密開集合の全閉点で、中心四変数上の階数256の自由表示と標準Koszul左右AS性が成立する。
\end{theorem}
\begin{proof}
前節の二つのグラフ行列を$\mathbb Q(b,c,v,w)$上で用いる。
\path{computations/verify_klein_generic_quartics.py}は四変数を代数的に独立のまま保ち、
各次数の関係行に共通する分母を払い、行の多項式最大公約数を除いて掃き出す。
各始点で次数0から5の次元は$1,4,10,20,35,56$である。
次数4から5への中心方程式は文字次数0で36列・階数33、次数1で40列・階数38となる。
分母は有理関数体の零でない元だけで、階数は有限体や数値近似から推定していない。
残る二文字次数は前出の有理点で核が零であり、そこで次数5までの商次元も一般値なので
一般点でやはり零である。各文字次数の核次元が有理点でも一般値となるため、
その点の近くで五次までを局所自由化すれば中心元の核も部分束になる。
以前に選んだ四式をこの部分束の局所切断に延ばし、
命題\ref{prop:klein-central-strata}の有限個の商次元上界が保たれる開集合を取ればよい。
この議論は基底の規格化に依存せず、行列補完の既約閉包の定階数部分に適用できる。
\end{proof}

\begin{corollary}\label{cor:klein-pi-specializations}
第二補完の閉包の任意の標準ASファイバーで、四頂点環は次数512の標準多項式恒等式
$s_{512}=0$を満たす。従ってこの閉包には非PIの標準AS窓はない。
\end{corollary}
\begin{proof}
上の開集合で四つの中心四次式の多項式環を$P$、四頂点環を$B$とする。
自由基底の次数多項式$(1+t+t^2+t^3)^4$では、指数を4で割った各剰余類の係数和は64である。
中心四変数が位置を変えないため$e_i B e_j$は全て$P$上自由階数64、
$B e_j$は自由階数256である。四つの$B e_j$上の左作用を合わせると忠実なので、
$B$は四つの$256$次行列環の直積に埋め込まれる。次数$2n$の標準恒等式を$n=256$に適用する。実際、外部変数$u_i$に対し
$T=\sum_{i=1}^{2n}u_i A_i$と置くと、奇元の巡回符号で$\operatorname{tr}(T^{2j})=0$である。
偶部分は可換なのでNewton恒等式とCayley--Hamiltonから$(T^2)^n=0$となり、
$u_1\cdots u_{2n}$の係数が所望の交代多項式である。
族の任意の標準AS点で任意の有限な内部次数まで、経路商の次元は一般値に等しい。
従ってその次数範囲では関係挿入行列の像は局所的に部分束である。
$s_{512}$の変数に固定した斉次な語を代入して得る有限次数の切断は、
一般点で零なのでこの近傍でも零となる。多重線型性で任意の代入を得る。
\end{proof}
